\section{Equations over unrestricted integers}
\label{sec:integers}

The distinction between positive and unrestricted integer unknowns is
substantial here.  Replacing all twelve positive unknowns by sums of
squares would discard much of the improvement.  Instead, we combine the
new compression with the signed-domain theorem of package 14.

\begin{lemma}[A three-variable positivity surjection]
The polynomial
\[
\Pos(x,y,z)=1+x^2+x+y^2+z^2
\]
has image exactly $\Z_{>0}$ on $\Z^3$.
\end{lemma}
\begin{proof}
The product $x(x+1)$ is nonnegative for integer $x$, so every value is
positive.  For $n>0$, Legendre's three-square theorem represents
$4n-3$ as three integer squares.  Since $4n-3\equiv1\pmod4$, exactly
one of the squares is odd.  Writing that coordinate as $2x+1$ and the
others as $2y,2z$ gives $n=1+x^2+x+y^2+z^2$.
\end{proof}

One copy costs seven binary operations.  The signed compiler theorem
allows us to retain only seven positive guards,
\[
J,F,\alpha,z,s,\eta,\zeta>0,
\]
provided $i$ is odd and $w$ is a square.  The other original witnesses
may be signed.  The canonical completeness proof supplies these two
restrictions; they are not imposed on arbitrary old solutions.

\subsection{Twenty-six unrestricted unknowns}

Let $h,T,t,r_0,u_0$ range over $\Z$, put
\[
i=2r_0+1,\qquad w=u_0^2,
\]
and replace each of the seven positive guards by its own independent
copy of $\Pos$.  This uses $7\cdot3+5=26$ unknowns.
Keep the arithmetic definitions of Section~\ref{sec:primitive} through
$A,D,Q,d,v,Z,b,G,e,B_*$.  The old separating weights cannot simply be
reused: $T$ is signed, and $w=0$ must still be permitted before the main
root is recovered.  Replace them by
\begin{align}
g&=c+X,&u&=\ell+k,&a_\tau&=2g+1,\nonumber\\
S&=a_\tau u+g,&C_*&=v^2+D+Z,&a_y&=QC_*S+1,\label{eq:integerweights}\\
r_f&=(a+2)(i^2+1)c^2,&C_2&=c(T^2+1)+c+Rr_f,
&a_f&=2a_yC_2+S+1.\nonumber
\end{align}
In the norm recipe use
\[
\mathcal Z_{26}=Ht+b-Ha_f\xi-Ha_\tau\theta
\]
and replace $\lambda$ by $a_y$.  Call the resulting polynomial
$\mathcal R_{26}$.  After the substitutions just specified, define
\begin{equation}
\boxed{\mathcal P_{26}=N_kN_t(1+\mathcal R_{26})-1.}
\label{eq:P26}
\end{equation}

\begin{theorem}
ZFC proves that $\mathcal P_{26}$ has no solution in $\Z^{26}$ if and
only if $\Con(\mathrm{ZFC})$.  Its exact degree is 5265 and its circuit
has 229 binary operations.  The corresponding three-equation system
has maximum degree 5248 and 225 operations.
\end{theorem}
\begin{proof}
First establish the bounds without assuming $X>0$.
The unit transport equation gives $C\ge0$ just as before, since
$K+w>1$, and hence $R>0$.  Although $X$ and $\ell$ may be zero,
$Y,k,c,a,H,\Delta,D,Q,d$ are positive, with $c>1$ and $D,Q>1$.
For each of the sixteen branches in~\eqref{eq:branches},
\begin{gather*}
|\gamma|<g,\qquad 1\le|\tau|<u,\qquad \sqrt D<r_f,\\
\frac1{QC_*}<|y|<2C_2\sqrt D.
\end{gather*}
The bound for $\gamma$ is unchanged.  The inequality
$u^2-A=\ell k+k^2-1>0$ also holds when $\ell=0$.
The chosen $r_f$ exceeds $\sqrt{1+\Delta i^2c^4}$ for every nonzero
integer $i$.  In the auxiliary estimates, replace $v$ by $|v|$.
We still have $|v|\sqrt D+Z<C_*$ and
$|v|+Z/\sqrt D<C_2$.

The product $e^2-B_*^2D$ cannot vanish.  When $B_*\ne0$ the earlier
argument applies.  When $B_*=0$, $e=QZ^2-1$ is a nonzero integer.
The conjugate-product estimate gives the lower bound for $|y|$.
For the upper bound, use $Q\ge4$, $Z>1$, and
\[
|y|^2\le\frac{Q(|v|\sqrt D+Z)^2+1}{Q-1}
 <4(|v|\sqrt D+Z)^2.
\]
These arguments allow signed $T$ and $i$ and do not involve $h$.

The sixteen sums
\[
\gamma+a_\tau\tau+a_y y+a_f f
\]
are distinct by the same successive-dominance argument.  In the
$\tau$ step, use $a_\tau|\tau|\ge a_\tau>g$; the stronger old
inequality $\sqrt A>g$ is no longer necessary.  The factorization has
the same nonzero scale $d^8H^{16}$.  A rational value $t$ therefore
makes all four coordinates rational.  The integrality arguments of
Section~\ref{sec:primitive} apply, using
$a_0=|\Delta i c^2|$ in Lemma~\ref{lem:rigidity}.

Only now exclude $w=0$.  The restored main equation would then give
\[
(ac+H\gamma)^2-(a^2+H)c^2=1,
\]
whose reduction modulo $H>1$ is $0\equiv1$.  Thus $w>0$.
The full signed theorem of package 14 now applies, with seven positive
guards, odd $i$, square $w$, and all reconstructed coordinates integral.
It gives a blank K215 halt and forces the remaining outer signs
positive.  The norm is a square modulo four by the same two-norm
argument, so the product finalizer is valid.

Conversely take the canonical signed-domain witness supplied by package
14, with the new coherent coefficient tuple.  Set
$t=\gamma+a_\tau\tau+a_y y+a_f f$.  This is an unrestricted integer,
so it needs no sign or parity restriction.  Surjectivity of $\Pos$
supplies the seven guard triples.  The resulting tuple solves
\eqref{eq:P26}.
\end{proof}

The ordering of this argument matters.  Invoking $X>0$ in the
separation estimates and later obtaining $X>0$ from a recovered main
root would be circular.  The new $a_\tau$ eliminates precisely that
problem.

After the quadratic guard substitutions, the degrees needed for the
leading-term calculation are
\begin{gather*}
\deg g=10,\quad \deg S=32,\quad \deg C_*=74,\quad\deg a_y=196,\\
\deg a_f=238,\quad\deg D=66,\quad\deg Q=90.
\end{gather*}
The initial $\mathcal K$ has a unique degree-656 leading contribution;
its competing contribution has degree 654.  Successive squaring and
norms give degrees $1312,2624,5248$.  At the final norm, the competing
term has degree 5184.  There is therefore no highest-degree
cancellation.  Since $\deg N_k=13$ and $\deg N_t=4$, the product has
degree $13+4+5248=5265$.  The sum of squares has degree 10496.

\subsection{Twenty-seven integers with a substantially lower degree}

A different combination preserves the main quotient $\gamma$ and
eliminates only $f,\tau,y$.  Keep the same seven positive guards, let
$h,T,\gamma,r_0,u_0,v_0$ be arbitrary integers, and set
\[
i=2r_0+1,\qquad w=u_0^2,\qquad t=2v_0+1.
\]
After the guard substitutions there are $7\cdot3+6=27$ unknowns.
With the earlier root data put
\begin{align}
N_m&=(X+ac+\gamma H)^2-\Delta c^2,\nonumber\\
M&=2(D+v^2),\qquad C_a=d(t^2+M^2A)-e,\nonumber\\
U&=C_a^2+DB_*^2-4d^2t^2M^2A,\label{eq:integer27recipe}\\
\mathcal S&=U^2-D(2C_aB_*)^2.\nonumber
\end{align}
The proposed polynomial is
\begin{equation}
\boxed{\mathcal P_{27}=N_mN_kN_t(1+\mathcal S)-1.}
\label{eq:P27}
\end{equation}
It combines the weighted amalgamation of package 09 with the signed
domain of package 14.  Doubling $M$ ensures the new compressed
coordinate can be odd in canonical witnesses.

At a zero, $N_m$ cannot be $-1$, since
$\Delta\equiv0$ or $3\pmod4$.  Similarly $\mathcal S$ is a square
modulo four.  Thus $N_m=1$, $\mathcal S=0$, and
$N_k=N_t=\varepsilon\in\{1,-1\}$.  The main norm excludes $w=0$.
The weak bounds then give $E>R+2$ and
$k=R+\varepsilon+hE>0$, so $h\ge0$ and $R\le k+1\le2k$.
Also $Y\ge4096$ and $c<k(Y+1)$.

Two possible zero denominators in the proof need explicit treatment.
First, $A>k^2Y^4>8(c^2+R^2)$, and
\[
(D+v^2)^2\ge4Dv^2,\qquad
Z^2\le2D^2(c^2+R^2).
\]
These imply $AM^2>2(Dv^2+Z^2)$.  Since $d\ge Q/2$, it follows
that $C_a>0$, including for signed $T$ and $h$.

Second, $B_*=0$ would mean $T=0$.  Then $\mathcal S=0$ gives
$C_a^2=4d^2t^2M^2A$.  Since $t$ is odd and nonzero,
$\tau=C_a/(2dtM)$ is rational with integer square $A$, and is therefore
integral.  The integer $y=t-M\tau$ satisfies
$dy^2=QZ^2-1$.  But $1<Z<Q$: use $Q=\Delta(D-1)$,
$\Delta>R$, and $D\ge1+\Delta c^4$ to get
\[
Q-Z=(\Delta-R)D-\Delta-c\ge D-\Delta-c>0.
\]
Consequently
\[
0<y^2-Z^2=\frac{Z^2-1}{Q-1}<Z,
\]
whereas the left side, as a positive difference of integer squares,
is at least $2Z+1$.  This contradiction excludes $B_*=0$.

All proof-only denominators in
\[
f=-\frac{U}{2C_aB_*},\qquad
\tau=\frac{C_a+B_*f}{2dtM},\qquad y=t-M\tau
\]
are now nonzero.  Direct substitution yields $f^2=D$, $\tau^2=A$,
and $dy^2=e-B_*f=Q(vf-Z)^2-1$.  The two rational square roots are
integral, and then $y$ is integral.  Package 14 proves soundness.
For completeness, its canonical $c,R$ and auxiliary ordinate $y$ are
odd.  Since $M$ is even, $t=M\tau+y$ is odd, so the new restriction
on $t$ is available.  Thus~\eqref{eq:P27} has exactly the required
existence predicate.

The arithmetic has 139 operations, $74M+65A$.  Its four factors have
degrees $35,13,4,1064$, so the exact product degree is 1116.
The degree-44 terms of $N_m$ cancel; its nonzero leading contribution
is $2ac\gamma H$ of degree 35.  In $\mathcal S$, the successive
unique top contributions are $dM^2A$, $C_a^2$, and $U^2$, of degrees
$266,532,1064$.  This is a lower-degree alternative to the
26-integer equation, rather than a smaller-variable replacement.

\begin{center}\small
\begin{tabular}{lrrrr}
\toprule
Unrestricted-integer construction & Variables & Equations & Degree & Operations\\
\midrule
Existing signed product (14) &29&1&188&121\\
New even-weight product &27&1&1116&139\\
New primitive-element system &26&3&5248&225\\
New primitive-element product &26&1&5265&229\\
New primitive-element sum of squares &26&1&10496&232\\
\bottomrule
\end{tabular}
\end{center}

These are improvements within the specified compiler family.
General eleven-integer representation results also exist
\cite{BayerDavid2025,BayerEtAl2025}; no global variable-count record is
claimed.  The attraction here is the much smaller displayed degree,
finite arithmetic circuit, and concrete specialization to the supplied
machine.
